% GENERATED FILE. Edit canonical lessons, then run npm run generate:books.
% canonical-source-hash: fnv1a64:3cb56ab5dd4a866a
% canonical-lessons: GU-C382-dhyey, GU-C382-sankalp, GU-C382-siddhi, GU-C382-ayojan, GU-C382-mathvu

\chapter{Setting a Goal}
\label{ch:gu-setting-a-goal}
\hlchaptermodality{382}

\begin{chapteropening}
\textbf{By the end of this chapter:} \emph{I can say what I want to reach and what it will take.}

Complete the last lesson of chapter 382: I can say what I want to reach and what it will take.
\end{chapteropening}

\section[\texorpdfstring{\gu{ધ્યેય} (dhyey)}{ધ્યેય (dhyey)}]{\gu{ધ્યેય} (dhyey) — a goal}

\label{lesson:GU-C382-dhyey}



\begin{quote}
Before the new one: say the Gujarati for to long for someone, then the Gujarati for an opening.
\end{quote}

\subsection*{You'll want to know: \gu{ધ્યેય}}
\textbf{\gu{ધ્યેય}} — \emph{dhyey} — \textquotedblleft{}a goal\textquotedblright{}.

\begin{quote}
\gu{જો} \gu{ધ્યેય} \gu{નક્કી} \gu{હશે} \gu{તો} \gu{મહેનત} \gu{કરવી} \gu{સહેલી} \gu{લાગશે}.
\end{quote}

\begin{quote}
\emph{jo dhyey nakkī haśe to mahenat karvī sahelī lāgśe.}
\end{quote}

\begin{quote}
If the goal is clear, the hard work will feel easier.
\end{quote}

A neuter noun.

\begin{cousinweb}
\textbf{\gu{ધ્યેય}} (\emph{dhyey}) is a tatsama, borrowed from Sanskrit \emph{dhyeya}, \textquotedblleft{}to be meditated on, to be pondered\textquotedblright{} (Monier-Williams): what the mind is fixed on.
\end{cousinweb}

\begin{grammarlens}[title={What you've built}]
One more word for saying what happens if.
\end{grammarlens}

\subsection*{Your turn}
\begin{itemize}
\raggedright
  \item \emph{Say it:} \emph{dhyey}
  \item \emph{Say it:} \emph{dhyey}, once more
  \item \emph{Say it:} \emph{jo dhyey nakkī haśe to mahenat karvī sahelī lāgśe.} — If the goal is clear, the hard work will feel easier.
  \item \emph{Recall:} say the Gujarati for to long for someone, then the Gujarati for an opening, then say \emph{dhyey} again
\end{itemize}

\begin{tcolorbox}[breakable,colback=teal!4,colframe=teal!35!black,title={Before you move on}]
What does \gu{ધ્યેય} mean? (\textquotedblleft{}A goal\textquotedblright{}.) Say it once more.
\end{tcolorbox}
\section[\texorpdfstring{\gu{સંકલ્પ} (sankalp)}{સંકલ્પ (sankalp)}]{\gu{સંકલ્પ} (sankalp) — a resolution}

\label{lesson:GU-C382-sankalp}



\begin{quote}
Before the new one: say the Gujarati for an opening, then the Gujarati for a goal.
\end{quote}

\subsection*{You'll want to know: \gu{સંકલ્પ}}
\textbf{\gu{સંકલ્પ}} — \emph{sankalp} — \textquotedblleft{}a resolution\textquotedblright{}.

\begin{quote}
\gu{જો} \gu{તમારો} \gu{સંકલ્પ} \gu{મજબૂત} \gu{હશે} \gu{તો} \gu{તમે} \gu{જૂની} \gu{ટેવ} \gu{છોડી} \gu{શકશો}.
\end{quote}

\begin{quote}
\emph{jo tamāro sankalp majbūt haśe to tame jūnī ṭev choḍī śakśo.}
\end{quote}

\begin{quote}
If your resolve is strong, you will be able to give up an old habit.
\end{quote}

A masculine noun.

\begin{cousinweb}
\textbf{\gu{સંકલ્પ}} (\emph{sankalp}) is a tatsama, borrowed from Sanskrit \emph{saṃkalpa}, \textquotedblleft{}a conception of the mind, a resolve\textquotedblright{} (Monier-Williams).
\end{cousinweb}

\begin{grammarlens}[title={What you've built}]
One more word for saying what happens if.
\end{grammarlens}

\subsection*{Your turn}
\begin{itemize}
\raggedright
  \item \emph{Say it:} \emph{sankalp}
  \item \emph{Say it:} \emph{sankalp}, once more
  \item \emph{Say it:} \emph{jo tamāro sankalp majbūt haśe to tame jūnī ṭev choḍī śakśo.} — If your resolve is strong, you will be able to give up an old habit.
  \item \emph{Recall:} say the Gujarati for an opening, then the Gujarati for a goal, then say \emph{sankalp} again
\end{itemize}

\begin{tcolorbox}[breakable,colback=teal!4,colframe=teal!35!black,title={Before you move on}]
What does \gu{સંકલ્પ} mean? (\textquotedblleft{}A resolution\textquotedblright{}.) Say it once more.
\end{tcolorbox}
\section[\texorpdfstring{\gu{સિદ્ધિ} (siddhi)}{સિદ્ધિ (siddhi)}]{\gu{સિદ્ધિ} (siddhi) — an achievement}

\label{lesson:GU-C382-siddhi}



\begin{quote}
Before the new one: say the Gujarati for a goal, then the Gujarati for a resolution.
\end{quote}

\subsection*{You'll want to know: \gu{સિદ્ધિ}}
\textbf{\gu{સિદ્ધિ}} — \emph{siddhi} — \textquotedblleft{}an achievement\textquotedblright{}.

\begin{quote}
\gu{જો} \gu{તે} \gu{આ} \gu{હરીફાઈ} \gu{જીતશે} \gu{તો} \gu{તે} \gu{તેની} \gu{સૌથી} \gu{મોટી} \gu{સિદ્ધિ} \gu{હશે}.
\end{quote}

\begin{quote}
\emph{jo te ā harīphāī jītśe to te tenī sauthī moṭī siddhi haśe.}
\end{quote}

\begin{quote}
If she wins this contest, it will be her greatest achievement.
\end{quote}

A feminine noun.

\begin{cousinweb}
\textbf{\gu{સિદ્ધિ}} (\emph{siddhi}) is a tatsama, borrowed from Sanskrit \emph{siddhi}, \textquotedblleft{}accomplishment, success\textquotedblright{} (Monier-Williams).
\end{cousinweb}

\begin{grammarlens}[title={What you've built}]
One more word for saying what happens if.
\end{grammarlens}

\subsection*{Your turn}
\begin{itemize}
\raggedright
  \item \emph{Say it:} \emph{siddhi}
  \item \emph{Say it:} \emph{siddhi}, once more
  \item \emph{Say it:} \emph{jo te ā harīphāī jītśe to te tenī sauthī moṭī siddhi haśe.} — If she wins this contest, it will be her greatest achievement.
  \item \emph{Recall:} say the Gujarati for a goal, then the Gujarati for a resolution, then say \emph{siddhi} again
\end{itemize}

\begin{tcolorbox}[breakable,colback=teal!4,colframe=teal!35!black,title={Before you move on}]
What does \gu{સિદ્ધિ} mean? (\textquotedblleft{}An achievement\textquotedblright{}.) Say it once more.
\end{tcolorbox}
\section[\texorpdfstring{\gu{આયોજન} (āyojan)}{આયોજન (āyojan)}]{\gu{આયોજન} (āyojan) — planning}

\label{lesson:GU-C382-ayojan}



\begin{quote}
Before the new one: say the Gujarati for a resolution, then the Gujarati for an achievement.
\end{quote}

\subsection*{You'll want to know: \gu{આયોજન}}
\textbf{\gu{આયોજન}} — \emph{āyojan} — \textquotedblleft{}planning\textquotedblright{}.

\begin{quote}
\gu{જો} \gu{આયોજન} \gu{સારું} \gu{હશે} \gu{તો} \gu{કાર્યક્રમ} \gu{સમયસર} \gu{પતશે}.
\end{quote}

\begin{quote}
\emph{jo āyojan sārũ haśe to kāryakram samaysar patśe.}
\end{quote}

\begin{quote}
If the planning is good, the programme will finish on time.
\end{quote}

A neuter noun.

\begin{cousinweb}
\textbf{\gu{આયોજન}} (\emph{āyojan}) is a tatsama, borrowed from Sanskrit \emph{āyojana}, \textquotedblleft{}junction, combination\textquotedblright{} (Monier-Williams); \textquotedblleft{}planning\textquotedblright{}, putting the parts together, is the Gujarati sense.
\end{cousinweb}

\begin{grammarlens}[title={What you've built}]
One more word for saying what happens if.
\end{grammarlens}

\subsection*{Your turn}
\begin{itemize}
\raggedright
  \item \emph{Say it:} \emph{āyojan}
  \item \emph{Say it:} \emph{āyojan}, once more
  \item \emph{Say it:} \emph{jo āyojan sārũ haśe to kāryakram samaysar patśe.} — If the planning is good, the programme will finish on time.
  \item \emph{Recall:} say the Gujarati for a resolution, then the Gujarati for an achievement, then say \emph{āyojan} again
\end{itemize}

\begin{tcolorbox}[breakable,colback=teal!4,colframe=teal!35!black,title={Before you move on}]
What does \gu{આયોજન} mean? (\textquotedblleft{}Planning\textquotedblright{}.) Say it once more.
\end{tcolorbox}
\section[\texorpdfstring{\gu{મથવું} (mathvũ)}{મથવું (mathvũ)}]{\gu{મથવું} (mathvũ) — to strive}

\label{lesson:GU-C382-mathvu}



\begin{quote}
Before the new one: say the Gujarati for an achievement, then the Gujarati for planning.
\end{quote}

\subsection*{You'll want to know: \gu{મથવું}}
\textbf{\gu{મથવું}} — \emph{mathvũ} — \textquotedblleft{}to strive\textquotedblright{}.

\begin{quote}
\gu{જો} \gu{તમે} \gu{સતત} \gu{મથશો} \gu{તો} \gu{ધ્યેય} \gu{સુધી} \gu{પહોંચશો}.
\end{quote}

\begin{quote}
\emph{jo tame satat mathśo to dhyey sudhī pahõcśo.}
\end{quote}

\begin{quote}
If you keep striving, you will reach your goal.
\end{quote}

Past (he …): \gu{તે} \gu{આખું} \gu{વરસ} \gu{મથ્યો}, \emph{te ākhũ varas mathyo}.

\begin{cousinweb}
\textbf{\gu{મથવું}} (\emph{mathvũ}) is inherited from Sanskrit \emph{mathnāti}, \textquotedblleft{}whirls, churns\textquotedblright{}; CDIAL lists the Gujarati verb, \textquotedblleft{}to churn\textquotedblright{} (CDIAL 9771). \textquotedblleft{}To strive\textquotedblright{}, to toil at a thing like churning milk, is a sense CDIAL does not give.
\end{cousinweb}

\begin{grammarlens}[title={What you've built}]
One more word for saying what happens if.
\end{grammarlens}

\subsection*{Your turn}
\begin{itemize}
\raggedright
  \item \emph{Say it:} \emph{mathvũ}
  \item \emph{Say it:} \emph{mathvũ}, once more
  \item \emph{Say it:} \emph{jo tame satat mathśo to dhyey sudhī pahõcśo.} — If you keep striving, you will reach your goal.
  \item \emph{Recall:} say the Gujarati for an achievement, then the Gujarati for planning, then say \emph{mathvũ} again
\end{itemize}

\begin{tcolorbox}[breakable,colback=teal!4,colframe=teal!35!black,title={Before you move on}]
What does \gu{મથવું} mean? (\textquotedblleft{}To strive\textquotedblright{}.) Say it once more.
\end{tcolorbox}
